\section{Elementary, algebraic, and polynomial special cases}\label{sec:special}

\subsection{The binomial and geometric series}

\begin{theorem}[Generalized binomial identity]\label{thm:binomial}
Formally,
\begin{equation}\label{eq:binomial}
 \F a b b z=\sum_{n\ge0}\frac{\poch a n}{n!}z^n=(1-z)^{-a}.
\end{equation}
For rational or algebraically named \(a\), the right side means the normalized
formal solution of
\((1-z)y'-ay=0\).  On a certified branch disk it is the named analytic power
\(\exp(-a\log(1-z))\).
\end{theorem}

\begin{proof}
The cancellation of \(\poch b n\) is finite and guarded.  If
\(B_n=\poch a n/n!\), then
\((n+1)B_{n+1}=(n+a)B_n\), which is exactly the coefficient recurrence of
\((1-z)y'-ay=0\), with \(B_0=1\).  The named power has the same differential
equation and initial value; effective uniqueness on a strict disk follows by
integrating the difference or by comparing Taylor coefficients and tails.
\end{proof}

\begin{corollary}[Geometric cases]
\begin{align*}
 \F{1}{b}{b}{z}&=\frac1{1-z}, &
 \F{1}{1}{1}{z}&=\frac1{1-z},\\
 \F{1}{2}{1}{z}&=\frac1{(1-z)^2}, &
 \F{1}{2}{2}{z}&=\frac1{1-z}.
\end{align*}
\end{corollary}

\begin{proof}
Apply \cref{thm:binomial}; for \(\F12 1z\), the coefficient is \(n+1\), and
the finite identity
\((1-z)^2\sum_{n<N}(n+1)z^n=1-(N+1)z^N+Nz^{N+1}\) supplies an explicit
remainder.
\end{proof}

\subsection{Logarithm and inverse sine}

\begin{theorem}[Logarithm]\label{thm:log-special}
On \(\norm z<1\), with the branch normalized by \(\log(1)=0\),
\[
 \F11 2{-z}=\frac{\log(1+z)}z,
\]
with the value at \(z=0\) interpreted as \(1\).
\end{theorem}

\begin{proof}
The coefficient is
\[
 \frac{\poch1n^2}{\poch2n n!}=\frac1{n+1}.
\]
Thus
\[
 z\F11 2{-z}=\sum_{m\ge1}(-1)^{m-1}\frac{z^m}{m}.
\]
Formal differentiation gives \((1+z)^{-1}\) and the constant term is zero.
Define the effective logarithm either by this series or by
\(\int_0^z(1+t)^{-1}\,dt\) on a star-shaped certified disk.  The geometric
tail proof in \cref{sec:convergence} makes the equality quantitative.
\end{proof}

\begin{theorem}[Arcsine]\label{thm:arcsin-special}
On a certified disk \(\norm z<1\),
\[
 \F{1/2}{1/2}{3/2}{z^2}=\frac{\arcsin z}{z},
\]
with value \(1\) at zero.
\end{theorem}

\begin{proof}
The finite factorial identities
\[
 \frac{\poch{1/2}n}{n!}=\frac1{4^n}\binom{2n}{n},
 \qquad
 \frac{\poch{1/2}n}{\poch{3/2}n}=\frac1{2n+1}
\]
are proved by induction after clearing powers of two.  Hence the coefficient
of \(z^{2n}\) is
\[
 \frac1{4^n(2n+1)}\binom{2n}{n}.
\]
Multiplying by \(z\) and differentiating gives the binomial series
\((1-z^2)^{-1/2}\), while the value at zero is zero.  This is the defining
effective integral for \(\arcsin z\) on the chosen branch.
\end{proof}

\subsection{The cubic radical identity}

The most striking elementary formula on the Wikipedia page can be proved
without guessing the radical: both sides are the uniquely normalized algebraic
solution of a cubic.

\begin{theorem}[Cubic algebraic hypergeometric function]\label{thm:cubic-radical}
Let
\[
 Y(x)=\F{1/3}{2/3}{3/2}{-27x^2/4}.
\]
As a formal power series,
\begin{equation}\label{eq:cubic-algebraic}
 Y+x^2Y^3=1,
 \qquad Y(0)=1.
\end{equation}
On the real branch with positive square and cube roots,
\begin{equation}\label{eq:cubic-radical}
Y(x)=
\frac{
 \sqrt[3]{\frac{3\sqrt3x+\sqrt{27x^2+4}}2}
 -\sqrt[3]{\frac{2}{3\sqrt3x+\sqrt{27x^2+4}}}
}{\sqrt3x},
\end{equation}
with the removable value \(1\) at \(x=0\).
\end{theorem}

\begin{proof}
Write \(t=x^2\).  Transforming the hypergeometric ODE gives
\begin{equation}\label{eq:cubic-ode}
 t(4+27t)Y''+(6+54t)Y'+6Y=0,
 \qquad Y(0)=1,\quad Y'(0)=-1.
\end{equation}
Now let \(U\) be the unique formal solution of \(U+tU^3=1\) with
\(U(0)=1\).  Existence and uniqueness are coefficientwise: the coefficient of
\(t^n\) in the cubic determines \(u_n\) from earlier coefficients.  Differentiate
the cubic twice, eliminate \(U'\) and \(U''\), and reduce with
\(tU^3=1-U\).  Exact polynomial simplification yields
\eqref{eq:cubic-ode}.  Its first two coefficients agree with those of \(Y\),
so the ODE recurrence gives \(U=Y\).  This entire calculation can be checked
in the finite truncated-series ring of \cref{def:trunc-series}.

For the radical, set
\[
 s=\sqrt{27x^2+4},\qquad
 u^3=\frac{s+3\sqrt3x}{2},\qquad v=u-u^{-1}.
\]
Because \((s+3\sqrt3x)(s-3\sqrt3x)=4\), one has
\(u^{-3}=(s-3\sqrt3x)/2\), and therefore
\[
 v^3+3v=u^3-u^{-3}=3\sqrt3x.
\]
Putting \(R=v/(\sqrt3x)\) gives \(R+x^2R^3=1\).  The positive isolated-root
choices make \(R\to1\) at zero, so formal/algebraic uniqueness gives
\(R=Y\).  Bisection names for the radicals and the strict-disk series modulus
turn the equality into an iPRA theorem schema.
\end{proof}

\subsection{A general algebraic family (Glasser's identity)}

\begin{theorem}[Well-poised algebraic family]\label{thm:glasser}
Let \(k\ge2\),
\[
 f_k=\frac{k}{(k+1)^{1+1/k}}.
\]
The numerator parameter vector
\(\frac1{k+1},\frac2{k+1},\ldots,\frac{k}{k+1}\), and denominator vector
\(\frac2k,\frac3k,\ldots,\frac{k-1}k,\frac{k+1}k\).  Then, as a formal
identity in \(x\),
\begin{equation}\label{eq:glasser}
 {}_kF_{k-1}\!\left(
 \begin{matrix}\frac1{k+1},\frac2{k+1},\ldots,\frac{k}{k+1}\\
 \frac2k,\frac3k,\ldots,\frac{k-1}k,\frac{k+1}k
 \end{matrix};
 \left(\frac{x(1-x^k)}{f_k}\right)^k\right)
 =\frac1{1-x^k}.
\end{equation}
For \(0\le x\le(k+1)^{-1/k}\), both sides have their ordinary positive-real
meaning and the convergence is effective.
\end{theorem}

\begin{proof}
Put \(u=x^k\),
\[
 t=\left(\frac{x(1-x^k)}{f_k}\right)^k,
 \qquad y=\frac1{1-u}.
\]
Since \(f_k^k=k^k/(k+1)^{k+1}\), direct algebra gives
\begin{equation}\label{eq:glasser-algebraic}
 y=1+f_k^k t\,y^{k+1}.
\end{equation}
The coefficient of the unique formal solution is the Fuss--Catalan number
\begin{equation}\label{eq:fuss}
 [t^n]y=\frac1{kn+1}\binom{(k+1)n}{n}f_k^{kn}.
\end{equation}
A fully finitary proof of \eqref{eq:fuss} is obtained by induction from
\eqref{eq:glasser-algebraic}: the coefficient of order \(n\) is a bounded
\((k+1)\)-fold convolution.  Alternatively, the standard Lagrange coefficient
formula is proved once by the finite binomial theorem and coefficient
extraction.

Apply the factorial multiplication identity
\[
 \prod_{j=1}^{k}\poch{j/(k+1)}n
 =\frac{((k+1)n)!}{(k+1)^{(k+1)n}n!}.
\]
After separating multiples of \(k+1\), and the analogous denominator
identity for the fractions with denominator \(k\).  Cross-multiplication
reduces the hypergeometric coefficient exactly to \eqref{eq:fuss}.  Hence the
left side is the unique formal solution of \eqref{eq:glasser-algebraic}, and
substitution gives \(y=(1-x^k)^{-1}\).  On the stated interval, the argument
lies in the convergence range and every root is the positive isolated root.
\end{proof}

\begin{corollary}[The \(k=2\) trigonometric-looking case]
For \(0\le x\le1/\sqrt3\),
\[
 \F{1/3}{2/3}{3/2}{\frac{27}{4}x^2(1-x^2)^2}
 =\frac1{1-x^2}.
\]
Combining this with \cref{thm:arcsin-special} or the cubic equation yields the
intermediate sine expression displayed by MathWorld; the algebraic equality
is the foundationally simpler statement.
\end{corollary}

\subsection{Jacobi, Legendre, and Chebyshev polynomials}

\begin{theorem}[Jacobi polynomial representation]\label{thm:jacobi}
For \(n\in\N\),
\begin{equation}\label{eq:jacobi}
 \F{-n}{n+\alpha+\beta+1}{\alpha+1}{x}
 =\frac{n!}{\poch{\alpha+1}n}P_n^{(\alpha,\beta)}(1-2x),
\end{equation}
where the right side may be characterized as the polynomial solution of the
Jacobi differential equation normalized by
\(P_n^{(\alpha,\beta)}(1)=\poch{\alpha+1}n/n!\).
\end{theorem}

\begin{proof}
The left side terminates at degree \(n\).  Substitute \(z=(1-t)/2\) into the
hypergeometric ODE.  Clearing denominators yields
\[
 (1-t^2)y''+[\beta-\alpha-(\alpha+\beta+2)t]y'
 +n(n+\alpha+\beta+1)y=0,
\]
the Jacobi equation.  The value at \(t=1\) is \(1\), so multiplication by
\(\poch{\alpha+1}n/n!\) gives the standard normalization.  Among polynomials
of degree at most \(n\), the coefficient recurrence determines the solution.
Everything is finite and remains valid over localized rational parameter
rings.
\end{proof}

\begin{corollary}[Legendre polynomials]
\[
 P_n(t)=\F{-n}{n+1}{1}{\frac{1-t}{2}}.
\]
\end{corollary}

\begin{theorem}[Generalized Chebyshev/cosine identity]\label{thm:chebyshev}
Formally in \(s\),
\[
 \F a{-a}{1/2}{\frac{1-s}{2}}=T_a(s).
\]
On a named real interval \(-\pi<x<\pi\), with \(s=\cos x\),
\[
 \F a{-a}{1/2}{\sin^2(x/2)}=\cos(ax).
\]
For integral \(a=n\), this is the ordinary Chebyshev polynomial identity and
is entirely finite.
\end{theorem}

\begin{proof}
Let \(z=\sin^2(x/2)\).  The chain rule transforms the hypergeometric equation
with parameters \((a,-a,1/2)\) into
\(Y''+a^2Y=0\).  The normalized solution has
\(Y(0)=1\), \(Y'(0)=0\), which characterizes \(\cos(ax)\) by its primitive
recursive Taylor coefficients and tail bound.  For integral \(n\), termination
makes \(T_n\) a polynomial and no trigonometric analysis is required.
\end{proof}

\begin{remark}[Legendre functions of nonintegral degree]
The source page also gives a formula involving an associated Legendre
function and gamma factors.  We treat it as a \StatusC{} identity: after a
specific Legendre branch is defined by its differential equation and local
normalization, the equality follows from the same change of variable and
formal uniqueness.  A global statement across its cuts additionally needs
branch-continuation data.
\end{remark}
